%% Original: ISCA 2022 Template
\documentclass[conference]{IEEEtran}
\usepackage{cite}
\usepackage{amsmath,amssymb,amsfonts}
\usepackage{algorithmic}
\usepackage{graphicx}
\usepackage{textcomp}
\usepackage{xcolor}
\usepackage[hyphens]{url}


\def\BibTeX{{\rm B\kern-.05em{\sc i\kern-.025em b}\kern-.08em
    T\kern-.1667em\lower.7ex\hbox{E}\kern-.125emX}}

% Ensure letter paper
\pdfpagewidth=8.5in
\pdfpageheight=11in

%%%%%%%%%%%%%%%%%%%%%%%%%%%%%%%%%%%%

\pagenumbering{arabic}

%%%%%%%%%%%---SETME-----%%%%%%%%%%%%%

\title{Paper Full Title}

\author{\normalsize{Conferenence Short Name (\textbf{Year})} \\ 
\normalsize{Name} \\ 
\normalsize{\today}}
%%%%%%%%%%%%%%%%%%%%%%%%%%%%%%%%%%%%

\begin{document}
\maketitle
\thispagestyle{plain}
\pagestyle{plain}


\section{Introduction}
What is the problem that is targeted by this paper and how it is a problem?
Main motivation of the paper and short description of the proposed solution. \\

Contributions of the paper \cite{nicepaper1}:
\begin{enumerate}
    \item Contribution-1
    \item Contribution-2
\end{enumerate}



\section{Proposal}
Main idea proposed. And specifics of the proposed solution. 
Should not contain low-level implementation details.

\section{Results}
Main results. Key takeaway points. Strengths of the paper.


\section{Limitations}

\subsection{Weaknesses}
Weaknesses or loopholes in the proposed idea. Any tradeoffs like more area for
performance. 

\subsection{Analysis Issues}
Statistical issues with the analysis. Improper baseline. Missing analysis.

\subsection{Unanswered Questions}
Some specifics of the idea not explained in the paper. Validity of the assumptions.


\section{Extensions}
Possible improvements on the idea. Different approach for the same problem.


%\section{References}

%%%%%%%%% -- BIB STYLE AND FILE -- %%%%%%%%
\bibliographystyle{IEEEtranS}
\bibliography{refs}
%%%%%%%%%%%%%%%%%%%%%%%%%%%%%%%%%%%%

%%%%%%% -- PAPER REVIEW ENDS -- %%%%%%%%

\end{document}
